% p09-inclusion-map-d-prime-into-d.tex
%
% Source:   Carl A. Gunter and Dana S. Scott, "Semantic Domains" (1990).
%           ScottDomains/papers/Gunter Scott 1990.pdf
% Physical page: 9           Printed page: 8
% Section:  2.3 "Uniformity", inside the proof of Theorem 3.
% Unnamed in the paper - no figure number, no caption.
%
% Introducing sentence (verbatim, printed page 8):
%   "Now, if i : D' \to D is the inclusion map then the following diagram
%    commutes"
% and the sentence closing it, immediately below the diagram:
%   "Thus, if F is a uniform fixed point operator, we must have
%    F_D(f) = F_{D'}(f')."
%
% What the development uses it for: it is the instance of the uniformity
% square (physical page 8) with h taken to be the inclusion of the sub-cpo
% D' = { x \in D | x \sqsubseteq fix(f) }, which is the step that forces any
% uniform fixed point operator to agree with fix.
%
% Geometry measured from a 150 dpi render of physical page 9: the two columns
% are 212 px apart and the two rows 120 px apart, i.e. 3.59 cm by 2.03 cm -
% the same square the paper prints on physical page 8, so the same 3.6 by 2.1
% is used here.
%
% Style contract (r0039) with the orchestrator-settled corrections, plus the
% "obj", "arrow" and "darrow" styles this commutative diagram needs (see
% p08-uniform-fixed-point-operator.tex for the note). All four arrows here
% are solid: the paper asserts no unique mediating map in this square.

\documentclass[border=6pt]{standalone}
\usepackage{tikz}
\begin{document}
\begin{tikzpicture}[
  x=1cm, y=1cm,
  every node/.style={inner sep=1pt},
  open/.style={circle, draw, fill=white, minimum size=4pt, inner sep=0pt},
  solid/.style={circle, draw, fill=black, minimum size=5pt, inner sep=0pt},
  edge/.style={draw, line width=0.4pt},
  obj/.style={inner sep=3pt},
  arrow/.style={draw, line width=0.4pt, ->},
  darrow/.style={draw, line width=0.4pt, ->, dash pattern=on 4pt off 3pt}]

  \node[obj] (Ptl) at (0.0, 2.1) {$D'$};
  \node[obj] (Ptr) at (3.6, 2.1) {$D'$};
  \node[obj] (Dbl) at (0.0, 0.0) {$D$};
  \node[obj] (Dbr) at (3.6, 0.0) {$D$};

  \draw[arrow] (Ptl) -- node[above] {$f'$} (Ptr);
  \draw[arrow] (Dbl) -- node[above] {$f$}  (Dbr);
  \draw[arrow] (Ptl) -- node[left]  {$i$}  (Dbl);
  \draw[arrow] (Ptr) -- node[right] {$i$}  (Dbr);

\end{tikzpicture}
\end{document}
